\customSection[sec:analysis]{Аналіз предметної області}{0em}

Перш ніж переходити до проєктування системи, доцільно розібратися в самій предметній області: що таке звичка з погляду поведінкової науки, чому її варто відстежувати й якими засобами це роблять сьогодні. Розуміння цих питань дає змогу обґрунтовано сформулювати вимоги та уникнути повторення вад, притаманних наявним рішенням.

\customSubSection[sec:analysis_domain]{Поняття звички та особливості її формування}{0em}

Звичка — це засвоєна модель поведінки, яка з часом починає виконуватися майже автоматично, без свідомого вольового зусилля. У дослідженнях поведінки людини усталилося уявлення про так звану «петлю звички», що складається з трьох елементів: сигналу (тригера), який запускає дію, самої дії та винагороди, що закріплює зв'язок між ними \citecustom{3}. Саме регулярне повторення цього циклу поступово перетворює свідому дію на автоматичну.

Окремо варто наголосити на ролі поступовості. Як показують праці, присвячені формуванню звичок, стійкого результату досягають не різкі й виснажливі зміни, а невеликі, але систематичні дії, які людина здатна повторювати щодня \citecustom{1}. Не менш важливим чинником є відчутність прогресу: коли людина бачить ланцюжок виконаних днів поспіль, у неї виникає додаткова мотивація не переривати його. Цей психологічний ефект, відомий як «не розривай ланцюг», лежить в основі більшості сучасних застосунків для трекінгу звичок \citecustom{2}.

З погляду обліку звички зручно поділити на два типи. \textbf{Бінарні} звички передбачають лише факт виконання («так» або «ні») — наприклад, «прочитати сторінку книги». \textbf{Вимірювані} звички пов'язані з кількісним показником і одиницею вимірювання — скажімо, «пробігти 2 кілометри» або «написати певну кількість рядків коду». Такий поділ безпосередньо впливає на структуру даних системи, адже для вимірюваних звичок потрібно зберігати не лише факт, а й конкретне значення.

Для фахівця у сфері розробки програмного забезпечення тема звичок має особливе значення. Професійне зростання тут безпосередньо залежить від регулярності: щоденне написання коду, систематичне читання документації, підтримання фізичної активності для збереження працездатності. Значна частина цієї діяльності вже залишає цифровий слід у зовнішніх сервісах, що відкриває можливість обліковувати її автоматично.

\customSubSection[sec:analysis_problem]{Проблематика ручного обліку}{1em}

Переважна більшість наявних застосунків для трекінгу звичок будується на ручному відмічанні виконаних дій. Користувач щодня заходить у застосунок і самостійно позначає, що звичку виконано. Попри простоту, такий підхід має низку суттєвих недоліків.

По-перше, він створює додаткове навантаження: щоб облік був корисним, людина має не забувати відмічати дії саме тоді, коли їх виконує. Парадоксально, але регулярне ведення обліку саме собою перетворюється на окрему звичку, яку так само важко сформувати й утримувати.

По-друге, страждає достовірність даних. Відмічання заднім числом, пропуски, забування призводять до того, що зібрана статистика лише приблизно відображає реальну поведінку. Ухвалювати рішення на основі таких даних складно.

По-третє, ручний облік ігнорує наявну цифрову активність користувача. Якщо розробник зробив десяток комітів у репозиторій, ця інформація вже зафіксована платформою GitHub, проте традиційний трекер усе одно вимагатиме окремо «відзвітувати» про виконану роботу. Виникає дублювання зусиль, якого можна було б уникнути.

\customSubSection[sec:analysis_analogs]{Огляд наявних рішень}{1em}

Щоб визначити місце розроблюваної системи, було розглянуто кілька популярних рішень як серед спеціалізованих трекерів звичок, так і серед сервісів, що фіксують активність користувача автоматично.

\textbf{Habitica} — гейміфікований трекер, який перетворює виконання звичок на рольову гру з персонажем, нагородами та штрафами. Сильний бік — мотивація через ігрові механіки, проте облік залишається повністю ручним, а інтеграцій з інструментами розробника немає.

\textbf{Streaks} та \textbf{Habitify} — мінімалістичні трекери для мобільних платформ, орієнтовані на підтримання щоденних ланцюжків. Streaks частково використовує дані Apple Health, однак ця автоматизація обмежена сферою здоров'я й замкнена в межах екосистеми Apple.

\textbf{Loop Habit Tracker} — відкритий застосунок для Android з гнучким налаштуванням розкладу та наочними графіками. Як і попередні, він покладається винятково на ручне введення.

\textbf{Граф внесків GitHub} (contribution graph) — не трекер звичок у строгому сенсі, а вбудована візуалізація активності розробника у вигляді теплової карти. Дані збираються автоматично, але прив'язані лише до подій GitHub і не дають змоги формулювати власні цілі чи звички.

\textbf{Strava} — соціальна платформа для спортсменів, яка автоматично реєструє тренування. Подібно до GitHub, вона фіксує активність без участі користувача, проте лише у власній, вузькій предметній області.

Зведене порівняння розглянутих рішень за ключовими критеріями наведено в \tabref{tab:analogs}.

\needspace{6cm}
\begin{spacing}{1.2}
    \footnotesize
    \noindent
    \begin{xltabular}{\textwidth}{| >{\raggedright\arraybackslash}X | c | c | c | c |}
        \caption{Порівняння наявних рішень за ключовими критеріями} \label{tab:analogs} \\
        \hline
        \textbf{Рішення} & \textbf{Веб-доступ} & \textbf{Автозбір} & \textbf{Свої звички} & \textbf{Інтеграції} \\ \hline
        \endfirsthead

        \hline
        \textbf{Рішення} & \textbf{Веб-доступ} & \textbf{Автозбір} & \textbf{Свої звички} & \textbf{Інтеграції} \\ \hline
        \endhead

        Habitica & частково & ні & так & ні \\ \hline
        Streaks & ні & частково & так & Apple Health \\ \hline
        Loop Habit Tracker & ні & ні & так & ні \\ \hline
        Граф внесків GitHub & так & так & ні & лише GitHub \\ \hline
        Strava & так & так & ні & лише спорт \\ \hline
        \textbf{Розроблювана система} & так & так & так & GitHub, Strava \\ \hline
    \end{xltabular}
\end{spacing}

\customSubSection[sec:analysis_feasibility]{Обґрунтування доцільності розробки}{1em}

Зіставлення розглянутих рішень виявляє чітку прогалину. Спеціалізовані трекери звичок (Habitica, Streaks, Loop) дають користувачеві свободу формулювати власні цілі, але покладаються на ручний облік. Натомість сервіси на кшталт графа внесків GitHub чи Strava збирають дані автоматично, проте замкнені у власній вузькій царині й не дають змоги визначати довільні звички. Жодне з розглянутих рішень не поєднує обидві переваги одночасно.

Саме на заповнення цієї прогалини й спрямована розроблювана система. Її ключова ідея полягає в тому, щоб дати користувачеві гнучкість звичайного трекера (довільні звички, бінарні та вимірювані, з власними цілями й періодичністю) і водночас зняти з нього тягар ручного обліку там, де активність уже фіксується зовнішніми сервісами. Коли користувач під'єднує GitHub або Strava, система самостійно опитує ці сервіси, виявляє нову активність і створює відповідні записи звичок без жодних дій з його боку. Там, де автоматизація неможлива, лишається класичне ручне введення.

Така комбінація робить розробку доцільною: вона усуває головну ваду наявних трекерів — залежність від дисциплінованості користувача в момент обліку — і підвищує достовірність зібраних даних, спираючись на об'єктивні події зовнішніх систем.

\customSubSection[sec:analysis_statement]{Постановка задачі та технічне завдання}{1em}

Спираючись на результати аналізу, задачу можна сформулювати так: необхідно створити вебзастосунок, який дає змогу користувачеві визначати власні звички, вести облік їх виконання та аналізувати прогрес, причому частина записів має формуватися автоматично на основі активності користувача в зовнішніх сервісах. Формально систему можна подати як перетворення множини подій — як ручних дій користувача, так і подій із сервісів GitHub та Strava — на впорядкований набір записів про виконання звичок з подальшим їх агрегуванням і візуалізацією.

\textbf{Вхідні дані:} облікові дані користувача (електронна пошта, пароль); параметри звичок (назва, тип, періодичність, ціль, теги, джерело автоматизації); токени доступу до зовнішніх сервісів; події активності, отримані із сервісів GitHub (коміти, запити на злиття) та Strava (тренування).

\textbf{Вихідні дані:} перелік звичок користувача та їх стан; записи про виконання звичок (ручні й автоматичні); зведена статистика та візуалізація активності у вигляді теплової карти й стрічки подій.

Перелік основних функцій системи та її акторів подано у вигляді діаграми використання (\textbf{рис.~\ref{fig:usecase}}). Головним актором є користувач, який взаємодіє з усіма основними сценаріями. Окремо виділено зовнішній сервіс як вторинного актора, що постачає дані для сценарію автоматичного збору активності.

\setfigure[fig:usecase]{0.78\textwidth}{assets/diagram-usecase}{Діаграма використання системи}

\textbf{Загальний опис інтерфейсу.} Взаємодія з користувачем відбувається через вебінтерфейс із бічним меню навігації, що містить розділи «Дашборд», «Звички», «Записи» та «Теги», а також окремий розділ профілю з налаштуваннями інтеграцій. Інтерфейс має бути адаптивним, із наочним поданням інформації та логічною структурою навігації.

Узагальнені технічні вимоги до системи наведено в \tabref{tab:tz}.

\needspace{7cm}
\begin{spacing}{1.2}
    \footnotesize
    \noindent
    \begin{xltabular}{\textwidth}{| >{\raggedright\arraybackslash}X | >{\raggedright\arraybackslash}X |}
        \caption{Технічне завдання (узагальнено)} \label{tab:tz} \\
        \hline
        \textbf{Характеристика} & \textbf{Вимога} \\ \hline
        \endfirsthead

        \hline
        \textbf{Характеристика} & \textbf{Вимога} \\ \hline
        \endhead

        Тип програмного продукту & вебзастосунок (односторінковий клієнт + REST API) \\ \hline
        Цільові користувачі & зареєстровані користувачі, насамперед розробники \\ \hline
        Автентифікація & обов'язкова, на основі токенів; захист реєстрації механізмом Altcha \\ \hline
        Зовнішні інтеграції & GitHub, Strava \\ \hline
        Періодичність автосинхронізації & кожні 2 хвилини \\ \hline
        Зберігання даних & реляційна СУБД; паролі й токени — у захищеному вигляді \\ \hline
        Доступ & через сучасний вебоглядач, без встановлення клієнта \\ \hline
    \end{xltabular}
\end{spacing}

\customSubSection[sec:analysis_requirements]{Вимоги до системи}{1em}

На основі проведеного аналізу сформульовано вимоги до системи, які поділено на функціональні та нефункціональні.

\customSubSubSection[sec:analysis_func]{Функціональні вимоги}{0em}

\begin{enumerate}[noitemsep, topsep=2pt, leftmargin=*, label=\arabic*)]
    \item реєстрація та автентифікація користувача із захистом від автоматизованих реєстрацій;
    \item створення, редагування та видалення звичок з можливістю задати тип (бінарна чи вимірювана), періодичність, цільове значення, кінцеву дату та проміжний рубіж (milestone);
    \item ведення записів про виконання звичок як вручну, так і автоматично;
    \item категоризація звичок за допомогою тегів;
    \item під'єднання зовнішніх сервісів (GitHub, Strava) до облікового запису користувача;
    \item автоматичне періодичне опитування під'єднаних сервісів та створення записів на основі виявленої активності;
    \item візуалізація активності користувача у вигляді теплової карти та стрічки останніх подій;
    \item перегляд і редагування профілю користувача.
\end{enumerate}

\customSubSubSection[sec:analysis_nonfunc]{Нефункціональні вимоги}{1em}

\begin{enumerate}[noitemsep, topsep=2pt, leftmargin=*, label=\arabic*)]
    \item \textbf{безпека} — зберігання паролів у захищеному вигляді, автентифікація на основі токенів, безпечне зберігання токенів доступу до зовнішніх сервісів;
    \item \textbf{розширюваність} — архітектура має дозволяти додавання нових інтеграцій та модулів без істотного переписування наявного коду;
    \item \textbf{підтримуваність} — чіткий поділ системи на шари та модулі з вираженими межами відповідальності;
    \item \textbf{надійність} — стійкість фонових процесів опитування до тимчасових збоїв зовнішніх сервісів;
    \item \textbf{зручність використання} — наочний та відгукливий вебінтерфейс;
    \item \textbf{спостережуваність} — можливість відстежувати роботу системи через структуроване журналювання й трасування.
\end{enumerate}

Сформульовані вимоги стають відправною точкою для проєктування архітектури системи, якому присвячено наступний розділ.
